\documentclass[10pt,a4paper]{amsart}
\usepackage[margin=19mm]{geometry}
\usepackage{amsmath,amssymb,mathtools}
\usepackage[scr=rsfs,cal=euler]{mathalfa}
\usepackage{xfrac}
\usepackage[unicode,hidelinks,bookmarks=false]{hyperref}
\newcommand{\ie}{i.e.\ }
\newcommand{\cf}{cf.\ }
\newcommand{\resp}{respectively\ }
\newcommand{\Subsection}{\subsection}
\providecommand{\nobreakdash}{\nobreak\hbox{-}}
\setlength{\emergencystretch}{2em}
\multlinegap0pt
\usepackage{amssymb}
\usepackage{mathtools}
\newtheorem{lemma}{Lemma}
\title{Partially Totally Real Submanifolds of Sasakian Manifolds}
\author{Chul Woo Lee, Jae Won Lee}
\date{Source: arXiv:2608.29406v1 (CC BY 4.0)}
\begin{document}
\maketitle

\noindent\emph{Source and licence.} Partially Totally Real Submanifolds of Sasakian Manifolds. Version arXiv:2608.29406v1 (CC BY 4.0), distributed by arXiv under the Creative Commons Attribution 4.0 International licence, http://creativecommons.org/licenses/by/4.0/, as recorded in the arXiv licence metadata for this exact version and shown on the abstract page https://arxiv.org/abs/2608.29406v1. Full licence notice: \texttt{licenses/arxiv-2608.29406-CC-BY-4.0.txt}.\par\smallskip

\noindent\textit{Editorial scope.} Counted unit: the article's lemma lem:Fparallel-shape (source0.tex lines 1084--1096). The article's complete original proof of it is delivered with it (source0.tex lines 1097--1136, one \texttt{proof} environment of 40 lines). No mathematics is written, repaired or re-derived: the statement and the proof are the article's own lines, in the article's own notation, and no retained line is substituted or re-flowed.  Retained uncounted as the source-local context the proof needs: lines 1030--1040.  Nothing was omitted from any retained span, so the package carries no omission record.  No retained line contains a citation, so the wrapper carries no bibliography and no bibliography entry.  No retained line contains a figure environment, an image-inclusion command or drawing code, so no image is delivered and the package declares no asset dependency.  Editorial formatting only, in addition: an AMS \texttt{amsart} preamble that restates the article's own declarations and macros used by the retained lines, namely the environments \texttt{lemma} on separate counters, together with the standard packages those lines need.  The retained environments print in this document as Lemma 1 in order of appearance, whereas the article prints the same items with the numbering of its own full document.  The complete native source of the article as distributed in the arXiv e-print for this exact version is bundled unchanged as \texttt{sources/208850/source0.tex}.
\begin{lemma}[Sasakian analogue of the K\"ahler $F$-formula]\label{lem:parallelF}
For all tangent vector fields $X,Y$,
\begin{equation}\label{eq:nablaF}
 (\bar\nabla_XF)Y=f h(X,Y)-h(X,PY).
\end{equation}
Hence $F$ is parallel if and only if
\begin{equation}\label{eq:Fparallel-h}
 h(X,PY)=f h(X,Y)
\end{equation}
for all $X,Y\in\Gamma(TM)$.
\end{lemma}

\begin{lemma}[Shape-operator characterization]\label{lem:Fparallel-shape}
The normal morphism \(F\) is parallel if and only if
\begin{equation}\label{eq:Fparallel-shape}
 A_{fV}X=P A_VX
\end{equation}
for every tangent vector field \(X\) and normal vector field \(V\).
Moreover, whenever \(F\) is parallel,
\begin{equation}\label{eq:Fparallel-anticommute}
 PA_V=-A_VP,
 \qquad
 A_{fV}=-A_VP .
\end{equation}
\end{lemma}

\begin{proof}
By Lemma~\ref{lem:parallelF}, \(F\) is parallel if and only if
\[
 h(X,PY)=f h(X,Y).
\]
Pair this identity with a normal vector field \(V\). Since \(\phi\) is
skew-adjoint and tangential and normal vectors are orthogonal,
\[
 g(fh(X,Y),V)
 =-g(h(X,Y),fV)
 =-g(A_{fV}X,Y),
\]
whereas
\[
 g(h(X,PY),V)
 =g(A_VX,PY)
 =-g(PA_VX,Y).
\]
Thus
\[
 g(A_{fV}X,Y)=g(PA_VX,Y)
\]
for every \(Y\), which is equivalent to
\(A_{fV}X=PA_VX\).  The calculation is reversible, so this condition is
also sufficient for \(F\) to be parallel.

Assume now that \(F\) is parallel.  Interchanging \(X\) and \(Y\) in
\(h(X,PY)=fh(X,Y)\) and using the symmetry of \(h\) gives
\[
 h(X,PY)=h(Y,PX).
\]
Pairing with \(V\) yields
\[
 g(A_VX,PY)=g(A_VY,PX).
\]
Using the skew-adjointness of \(P\) and the self-adjointness of \(A_V\),
this is equivalent to \(PA_V=-A_VP\). Combining this with
\(A_{fV}=PA_V\) gives the second identity in
\eqref{eq:Fparallel-anticommute}.
\end{proof}

\end{document}
